\chapter{Referencial teorico}
\label{cap:referencial}

O referencial mostra o que ja se sabe sobre o tema e onde o seu trabalho se
encaixa nessa conversa. Nao e um resumo de textos em sequencia: e a construcao
de um argumento apoiado neles.

\section{Primeiro conceito}
\label{sec:conceito-um}

Defina os conceitos que o trabalho usa, apoiando-se em quem os formulou
\cite{knuth1984}.

\section{Trabalhos relacionados}
\label{sec:relacionados}

Compare o que outros autores fizeram, apontando semelhancas, diferencas e
lacunas. A Tabela~\ref{tab:comparativo} ilustra uma forma de organizar essa
comparacao.

\begin{table}[htb]
  \centering
  \caption{Comparativo entre os trabalhos relacionados}
  \label{tab:comparativo}
  \begin{tabular}{lccc}
    \toprule
    \textbf{Trabalho} & \textbf{Abordagem} & \textbf{Escala} & \textbf{Avaliacao} \\
    \midrule
    \citeonline{exemplo2024} & Quantitativa & Media  & Experimental \\
    Este trabalho            & Mista        & Grande & Experimental \\
    \bottomrule
  \end{tabular}
  \fonte{Elaborado pelo autor (2026).}
\end{table}
